% !TEX TS-program = xelatex
% !TEX encoding = UTF-8 Unicode
% !Mode:: "TeX:UTF-8"

\documentclass{resume}
\usepackage{zh_CN-Adobefonts_external} % 简体中文字体支持（外部字体）
\usepackage{linespacing_fix}
\usepackage{cite}
\usepackage[hidelinks]{hyperref}
\usepackage{url}
\usepackage{tikz}
\usepackage{float}
\usepackage{wrapfig}
\hypersetup{
  colorlinks=true,
  linkcolor=black,
  urlcolor=black,
  filecolor=black,
  citecolor=black,
}
\usepackage{graphicx}  % 如果未添加，请加上此行
\begin{document}
\pagenumbering{gobble}


\name{李展利}
\contactInfo{\faPhone\ (+86)156-2350-1965}{\faEnvelope\ lizhanli@stu.zuel.edu.cn}
{\href{https://zhanli-li.github.io/}{\faHome\ zhanli-li.github.io}}
{\href{https://github.com/Zhanli-Li}{\faGithub\ Zhanli-Li (190 stars+)}}
%----------------------------------------------------------------------------------------
% 教育背景
%----------------------------------------------------------------------------------------
% 在 \name 命令之后添加
% \begin{tikzpicture}[remember picture,overlay]
%   \node[anchor=north east,xshift=-0.1cm,yshift=-0.35cm] at (current page.north east) {
%     \includegraphics[width=2.5cm,height=3.5cm,keepaspectratio]{头像.png}
%   };
% \end{tikzpicture}
\section{教育背景}
\datedsubsection{\textbf{中南财经政法大学} \quad \textbf{文澜学院(荣誉学院)} \quad \textit{数字经济}}{2023.9 - 2027.6}
\textbf{加权平均分:} 93.73/100 \quad \textbf{专业排名:} 2/80 \quad \textbf{奖学金:} \textbf{2025 年国家奖学金} \\
\textbf{主修课程:} 神经网络 \quad 机器学习 \quad 数学建模 \quad 数学分析 \quad 概率论与数理统计\\
\textbf{研究兴趣:} LLM \quad Agentic Training \quad Document Intelligence
\quad Causal Inference


%----------------------------------------------------------------------------------------
% 学术论文
%----------------------------------------------------------------------------------------
\section{学术论文}
\textbf{Zhanli Li}, Huiwen Tian, Lvzhou Luo, Yixuan Cao, Ping Luo. \textit{DeepRead: Document Structure-Aware Reasoning to Enhance Agentic Search}. \textit{KDD 2027} (CCF-A), \textbf{Resubmit} \quad \textbf{[First Author]}

\textbf{Zhanli Li}, Yixuan Cao, Lvzhou Luo, Ping Luo. \textit{Navigating Large-Scale Document Collections: MuDABench for Multi-Document Analytical QA}. \textit{ACL 2026 Findings} (CCF-A) \quad \textbf{[First Author]}
 
\textbf{Zhanli Li}, Zichao Yang. \textit{ESG Rating Disagreement and Corporate Total Factor Productivity: Inference and Prediction}. \textit{Finance Research Letters}, vol.~78, p.~107127, 2025. (中科院 Q1 Top，JCR Q1，IF: 6.9） \quad \textbf{[First Author]}

%----------------------------------------------------------------------------------------
% 项目经历
%----------------------------------------------------------------------------------------
\section{项目经历}
\datedsubsection{\textbf{Data Agent Training of LongCat 2.0\textsuperscript{‡}} \quad \textit{LongCat Interaction 研究实习生}}{2026.5 - 至今}
针对基座模型数据分析中\textbf{下钻不足、洞察停留表层}的问题，围绕 \textbf{LongCat 2.0 Pro 1.6T} 设计 \textbf{Data Agent} 后训练方案；构建\textbf{合成交互式分析环境生成 pipeline}、\textbf{轨迹筛选 pipeline}，并探索与其他 post-training 数据的\textbf{混合配比策略}，相关方案在 \textbf{LongCat 2.0 Pro 1.6T} 上完成验证，并适配\textbf{超万张国产卡（H910-64G）}训练体系，支撑其在 \textbf{BA Agent / Data Agent} 任务中的\textbf{下钻能力}与 \textbf{insight 质量}提升。\\
$\triangleright$~ 搭建了 Env 合成流程、深度分析轨迹筛选流程，以及 Data Agent 数据在后训练中的配比策略，并在 \textbf{LongCat 2.0 Pro 1.6T} 与\textbf{超万张国产卡（H910-64G）}训练体系上完成验证

\datedsubsection{\textbf{DeepRead：文档结构感知的 Agentic Search\textsuperscript{†}} \quad \textit{第一作者}}{2025.10 - 2026.2}
提出 \textbf{DeepRead} 结构感知文档推理 Agent，利用 OCR 保持版面结构保真，构建\textbf{段落级坐标}，并设计 \textbf{Retrieve} 与 \textbf{ReadSection} 工具，使模型形成 ``locate-then-read'' 推理范式；在四个多类型文档基准上较 \textbf{Search-o1} 基线平均提升 \textbf{10.3\%}。\\
$\triangleright$~ \textbf{Resubmit} \textit{KDD 2027 (CCF-A)}，预印本被知名媒体\textbf{新智元}报道

\datedsubsection{\textbf{MuDABench：多文档分析问答基准与多 Agent 框架\textsuperscript{†}} \quad \textit{第一作者}}{2025.4 - 2025.12}
构建 \textbf{MuDABench} 多文档分析问答基准（\textbf{80k+页}文档、\textbf{332}道问题、\textbf{4,964}条结构化中间事实），用于评测多文档检索、跨文档聚合与数值推理能力；设计 \textbf{Multi-Agent} 工作流，显著提升标准检索基线表现。\\
$\triangleright$~ \textbf{发表于} \textit{ACL 2026 (CCF-A)}

\datedsubsection{\textbf{基于 SHAP 的公司金融因果推断框架\textsuperscript{*}} \quad \textit{项目负责人}}{2024.9 - 2025.1}
构建可解释机器学习框架，结合编码面板数据与 \textbf{Optuna} 自动超参搜索，在 TFP 预测任务上取得 $R^2=0.76$；进一步使用 \textbf{SHAP} 量化 ESG 评级分歧对结果变量的\textbf{非线性边际贡献}，揭示潜在因果路径。\\
$\triangleright$~ \textbf{发表于} \textit{Finance Research Letters}（中科院 Q1 Top，JCR Q1，IF: 6.9），获清华大学因果推断研讨会优秀论文

{\small \textit{\textsuperscript{†}\textbf{中科院计算所}曹逸轩、罗平指导；\textsuperscript{‡}\textbf{LongCat Interaction} 洪峰指导；\textsuperscript{*}\textbf{文澜学院}杨子超指导}}

%----------------------------------------------------------------------------------------
% 实习经历
%----------------------------------------------------------------------------------------
\section{实习经历}
\datedsubsection{\textbf{美团 LongCat Interaction（基座大模型）} \quad \textit{研究实习生}}{2026.5 - 至今}
\textbf{导师:} 洪峰 \quad \textbf{项目:} BA Agent 能力提升；参与 \textbf{LongCat 2.0 Pro 1.6T} 在真实业务数据分析场景中的后训练与能力增强，并在\textbf{超万张国产卡（H910-64G）}训练体系中完成相关方案验证。

\datedsubsection{\textbf{北京庖丁科技有限公司（AI Startup）研究组} \quad \textit{AI技术研究实习生}}{2025.6 - 2026.2}
\textbf{导师:} 罗平、曹逸轩 \quad \textbf{项目:} 上下文检索在实际生产环境中的改进；基于 \textbf{PDFlux} 优化 \textbf{ChatDOC} RAG 系统上下文缺失问题。

%----------------------------------------------------------------------------------------
% 竞赛经历
%----------------------------------------------------------------------------------------
\section{竞赛经历}
\datedsubsection{\textbf{NVIDIA Nemotron Model Reasoning Challenge} \quad \textit{银牌}}{2026.6}
基于 \textbf{Nemotron-3-Nano-30B-A3B} 与比赛要求的 \textbf{LoRA} 微调设置，构建 Agentic 推理框架求解抽象谜题，并将归纳规则编排为可复用 solver，最终取得 \textbf{0.86} 得分，全球排名119/4182。

\section{技能与服务}
\textbf{编程/工具:} Python, C/C++, \LaTeX, Markdown, Docker, SSH, tmux, Ubuntu \quad \textbf{英语:} CET-4：522 \quad CET-6: 508\\
\textbf{框架:} ms-swift, vllm, verl, llamaindex, transformers, torch \quad



\end{document}
